\section{Conclusion}
The strongest result of this study is an ordering result rather than a detector result. When authorization constrains the retrieval candidate set before context construction, unauthorized records do not enter model context in any of 1,750 evaluated unauthorized requests; prompt-only policy leaves that boundary open in all 1,750. Generated canary disclosure is a distinct and much rarer outcome, and its observed 5-to-0 reduction is not overclaimed. A controlled retrieval-depth experiment further shows that the originally observed authorized-task gap was caused by distractor composition in this benchmark, not by the authorization gate itself. Finally, held-out paraphrases expose a clear limitation of heuristic risk triage: ACL denial is phrasing-invariant, while suspicion detection is not, and fuzzy aggregation shows no advantage over transparent rules. At scale, growing candidate pools increase retrieval cost and distractor similarity, reinforcing the systems case for reducing the search space before sensitive context reaches a model. Authorization-before-retrieval should therefore be treated as a structural boundary; risk scoring is a secondary policy layer, not a substitute for it.
